\section{The deterministic brane-lattice substrate}
\label{sec:substrate}

\subsection{Ontology and action}
\label{subsec:ontology}

The fundamental object is a four-dimensional cubic \emph{brane lattice} embedded with
codimension zero in a 4D Euclidean ambient $\R^4$. Lattice sites are labeled by
$n\in\Z^4$; each site carries an embedded position $R_n^A\in\R^4$ ($A=1,\dots,4$), and
the reference (undeformed) configuration is $R_n^A=a\,n^A$ with lattice spacing $a$. The
ambient is fully symmetric---no preferred direction---and serves only to measure Euclidean
distances. All asymmetries (the timelike direction; the gauge/gravity split) are
properties of the brane \emph{action}, not of the ambient.

The action is a single central-force spring sum over the $8$-link stencil (six spacelike
axial neighbours $\pm\hat{e}_i$, two temporal $\pm\hat{e}_4$):
\begin{equation}
  S[R] \;=\; \tfrac12 \sum_{n}\sum_{\mu=1}^{4}
      \eta_\mu\,\kappa_\mu \bigl(L_{n\mu}-r_\mu\bigr)^2,
  \qquad
  L_{n\mu}=\Bigl(\textstyle\sum_A (R_{n+\hat e_\mu}^A-R_n^A)^2\Bigr)^{1/2},
  \label{eq:action}
\end{equation}
with the Lorentzian sign pattern
\begin{equation}
  \eta_\mu=(-1,-1,-1,+1).
  \label{eq:eta}
\end{equation}
The link stiffnesses and stress-free (rest) lengths are grouped by direction,
\begin{equation}
  \kappa_i=\kappa_s,\quad \kappa_4=\kappa_t\;(=\gamma_t\kappa_s);\qquad
  r_i=\alpha_s a,\quad r_4=\alpha_t a \quad(i=1,2,3),
  \label{eq:params}
\end{equation}
so the minimal dimensionless control set is $\{\alpha_s,\alpha_t,\gamma_t\}$ with
$\gamma_t=\kappa_t/\kappa_s$, at fixed held spacing $a$. The single nonlinearity is the
Euclidean square root in $L_{n\mu}$; even a Hooke-linear scalar spring law is nonlinear in
the node coordinates because the extension is the full ambient distance. This is the
microscopic origin of the amplitude--lateral coupling that drives every emergent structure
below.

\paragraph{Signature as prestress sign, not postulate.} The timelike direction is selected
by $\eta_4=+1$ opposing $\eta_i=-1$: temporal link contributions enter the action with the
opposite sign from spacelike ones. The signed prestress is
$\rho_\mu = \eta_\mu\kappa_\mu a(1-\alpha_\mu)$; a nonzero mismatch $a-r_\mu$ under
periodic boundary conditions holds the links in tension. As shown in
Section~\ref{sec:metric}, this sign pattern is exactly the signature of the emergent
metric---Lorentzian signature is derived from prestress, not imposed alongside an ambient
metric. Prestress magnitude and sign are logically separate: $\alpha_\mu=1$ is the
zero-mismatch (no-prestress) limit, at which the transverse stiffness $\propto(1-\alpha_\mu)$
vanishes; the working regime keeps $0<\alpha_s\le 1$ and $0<\alpha_t<1$ (typically
$\alpha_t=1-\varepsilon_t$; see Section~\ref{sec:temporal}).

\paragraph{Block and initial-value views.} Stationarity $\partial S/\partial R_n^A=0$ is,
depending on boundary conditions, either a forward Störmer--Verlet march (initial-value
problem) or a 4D boundary-value ``block'' root-find. Because $S$ is indefinite
(Lorentzian, $S=T-V$), one root-finds $\nabla S=0$; one never minimizes. The two views
share the same local Euler--Lagrange stencil, which is precisely the variational (Verlet)
integrator \citep{MarsdenWest2001}. The parameterization~\eqref{eq:params} refines the
earlier single-$\alpha$, $r_t\approx 0$ formulation, in which the timelike sign could be
mistaken for a kinetic-sign artifact rather than a genuine prestress; the resolution is
Section~\ref{sec:temporal}.

\subsection{The anisotropic stiffness tensor}
\label{subsec:stiffness}

Let $\bar R$ be a stationary background ($\partial S/\partial R[\bar R]=0$) and write
fluctuations $R_n=\bar R_n+\xi_n$. The background link vectors and directions are
\begin{equation}
  \bar Q_{n\mu}^A=\bar R_{n+\hat e_\mu}^A-\bar R_n^A,\qquad
  \bar L_{n\mu}=|\bar Q_{n\mu}|,\qquad
  \Qhat_{n\mu}^A=\bar Q_{n\mu}^A/\bar L_{n\mu}.
\end{equation}
The exact second variation of a single link energy with respect to its link vector is the
\emph{anisotropic stiffness tensor}
\begin{equation}
  \boxed{\;
  C_{n\mu}^{AB}
  = \eta_\mu\kappa_\mu\!\left[
      \Bigl(1-\tfrac{r_\mu}{\bar L_{n\mu}}\Bigr)\delta^{AB}
      + \tfrac{r_\mu}{\bar L_{n\mu}}\,\Qhat_{n\mu}^A\Qhat_{n\mu}^B
    \right].
  \;}
  \label{eq:stiffness}
\end{equation}
It is an isotropic part plus a rank-one part along the link: a spring is stiffer along its
own axis than transverse to it. Its eigenvalues are $\eta_\mu\kappa_\mu(1-r_\mu/\bar L_{n\mu})$
(transverse, threefold) and $\eta_\mu\kappa_\mu$ (longitudinal); the transverse bracket is
positive whenever $r_\mu<\bar L_{n\mu}$ (stretched links), a fact we use for stability in
Section~\ref{sec:limits}.

In the trivial background $\bar Q_{n\mu}^A=a\,\delta_\mu^A$, so $\Qhat_{n\mu}^A=\delta_\mu^A$
and $C_{n\mu}$ is diagonal in the fixed axis frame; all $C_{n\mu}$ commute and share a
$k$-independent eigenbasis, giving a trivial eigenbundle. In a finite-amplitude nonlinear
background the $\Qhat_{n\mu}$ differ link to link, so generically
$[C_{n\mu},C_{m\nu}]\neq 0$. This link-to-link noncommutativity is the mechanical seed of
the non-abelian sector (Section~\ref{sec:carrier}); it is a property that lives \emph{in}
the Pythagorean nonlinearity, and vanishes if one linearizes it away. Consequently, the
absence of a non-abelian sector in the bare straight vacuum is not a no-go theorem---it
merely says the trivial background is the wrong place to look.

\subsection{The quadratic fluctuation action and the Bloch operator}
\label{subsec:bloch}

The quadratic fluctuation action is
\begin{equation}
  S^{(2)}[\xi;\bar R]
  = \tfrac12\sum_{n,\mu}\bigl(\xi_{n+\hat e_\mu}-\xi_n\bigr)^{\!A}
    C_{n\mu}^{AB}\bigl(\xi_{n+\hat e_\mu}-\xi_n\bigr)^{\!B}.
  \label{eq:s2}
\end{equation}
For a background periodic with a supercell (sites $s$, integer supercell offsets), the
Bloch ansatz $\xi_{n}=\varepsilon_s(\veck)\,\e^{\ii \veck\cdot n}$ turns each link into a
matrix-valued hopping $C_\ell^{AB}\e^{\ii\veck\cdot\delta}$ and yields a Hermitian Bloch
fluctuation operator $D_{\bar R}(\veck)$ with
\begin{equation}
  S^{(2)}=\tfrac12\sum_{\veck}\varepsilon^\dagger(\veck)\,D_{\bar R}(\veck)\,\varepsilon(\veck).
  \label{eq:Dk}
\end{equation}
The complex phase $\e^{\ii\veck\cdot n}$ is the character of the 4D lattice translation
group, used to diagonalize translations; physical configurations remain real through
conjugate pairs $\veck,-\veck$. Because $\eta_4=+1$ while $\eta_i=-1$, $D_{\bar R}(\veck)$
is Hermitian but \emph{indefinite} (the Lorentzian signature is inherited by the spectrum);
its eigenvalues are real and its spectral projectors well-defined. The operator
$D_{\bar R}(\veck)$ is the central object for both the emergent metric
(Section~\ref{sec:metric}, trivial background) and the gauge sector
(Section~\ref{sec:carrier}, nonlinear carrier).
